% =============================================================================
% Mambo — Closed-Form Proofs for Open Gaps
% Draft: t >= 3 structural lemma, multi-block CICO, indecomposability
% =============================================================================

\documentclass[a4paper,11pt]{article}
\usepackage{amsmath,amssymb,amsthm}
\usepackage{proof}

\newtheorem{lemma}{Lemma}
\newtheorem{proposition}{Proposition}
\newtheorem{theorem}{Theorem}
\newtheorem{corollary}{Corollary}
\newtheorem{definition}{Definition}

\newcommand{\Fp}{\mathbb{F}_p}
\newcommand{\Z}{\mathbb{Z}}
\newcommand{\GF}{\operatorname{GF}}
\newcommand{\disc}{\operatorname{disc}}
\newcommand{\lc}{\operatorname{lc}}
\newcommand{\LM}{\operatorname{LM}}
\newcommand{\inI}{\operatorname{in}(I)}
\newcommand{\R}{\mathcal{R}}
\newcommand{\B}{\mathcal{B}}

\begin{document}

\title{Mambo: Closing the Open Proof Gaps}
\author{Manuel Jordan Basso}
\date{}
\maketitle

%=============================================================================
\section{Structural Gr\"obner Basis for $t = 3$}\label{sec:t3}
%=============================================================================

We prove that for the rate-$1$ sponge CICO ideal at state width $t = 3$
and every $n \ge 2$, the quotient ring satisfies
\[
\dim_{\Fp} R/I = 2^{3n-2},\qquad
\operatorname{Hilb}_{R/I}(z) = (1+z)^{3n-2},\qquad
d_{\mathrm{reg}} = 3n-1.
\]

\subsection{The CICO Ideal}

Let $R = \Fp[x_{i,j} : 0 \le i \le n,\, 1 \le j \le 3]$. The Mambo-$3$
round equations (Section~\ref{sec:inst:wider}) are
\begin{align*}
x_{i+1,1} &= \varphi x_{i,1} + x_{i,3}^2 + c_{i,1},\\
x_{i+1,2} &= \varphi x_{i,2} + x_{i+1,1}^2 + c_{i,2},\\
x_{i+1,3} &= \varphi x_{i,3} + x_{i+1,2}^2 + c_{i,3},
\end{align*}
for $i = 0,\dots,n-1$, where $\varphi \in \Fp^\ast$ and
$c_{i,j}$ are distinct per-round constants. The sponge CICO constraints
fix the capacity at input and the first output coordinate:
\begin{align*}
x_{0,1} &= 0\quad (\text{capacity input}),\\
x_{0,2} &= 0\quad (\text{capacity input}),\\
x_{n,1} &= \tau\quad (\text{output constraint}).
\end{align*}
The rate cell $x_{0,3}$ is free.

Let $I \subseteq R$ be the ideal generated by the $3n+3$ polynomials above.

\subsection{Fixed Variables}

\begin{lemma}[Linear elimination]\label{lem:fixed_t3}
In the reduced degrevlex Gr\"obner basis of $I$, the following variables
have linear leading monomials:
\[
\mathcal{X}_n = \{x_{0,1},\, x_{0,2},\, x_{n-2,3},\, x_{n-1,2},\, x_{n,1}\}.
\]
Hence $|\mathcal{X}_n| = 5$ and these variables are eliminated from the
standard monomials.
\end{lemma}

\begin{proof}
The constraint $x_{0,1} = 0$ and $x_{0,2} = 0$ are generators, so their
LMs are themselves. The output constraint $x_{n,1} = \tau$ propagates
backward: from the round equations, $x_{n,1}$ is expressible in terms of
$x_{n-1,1}$ and $x_{n-1,3}^2$, and inductively
$x_{n-1,2}$ and $x_{n-2,3}$ appear linearly. The complete elimination
chain gives linear LMs for $x_{n-1,2}$ and $x_{n-2,3}$.

Explicitly, from the last round:
\[
x_{n,1} = \varphi x_{n-1,1} + x_{n-1,3}^2 + c_{n-1,1},
\]
and substituting $x_{n,1} = \tau$ gives a relation involving
$x_{n-1,1}$ and $x_{n-1,3}$. Then $x_{n-1,2}$ appears linearly in the
expression for $x_{n-1,3}$, and $x_{n-2,3}$ appears linearly in
$x_{n-1,2}$. The specific degrevlex ordering makes these linear terms
the leading monomials.
\end{proof}

\subsection{Free Variables and Their Constraints}

After Gauss-reducing the linear LMs, $3n+3 - 5 = 3n-2$ variables remain.
We partition them into two groups.

\begin{definition}[Interior free variables]\label{def:interior_t3}
The \emph{interior} free variables are
\[
\mathcal{I}_n = \mathcal{F}_n \setminus \{x_{n-1,3},\, x_{n,2},\, x_{n,3}\},
\]
where $\mathcal{F}_n$ is the complete free variable set
$\mathcal{F}_n = R_{\ge 1} \setminus \mathcal{X}_n$.
These variables occupy cells that are not part of the last cascade step
and have $i \le n-2$ or intermediate positions.
\end{definition}

\begin{lemma}[Square reduction for interior vars]\label{lem:square_t3}
For every $f \in \mathcal{I}_n$, the Gr\"obner basis contains a
generator with leading monomial $f^2$. Consequently every standard
monomial has $\deg f \le 1$.
\end{lemma}

\begin{proof}
For $x_{i,1}$ with $1 \le i \le n-1$: the equation
$x_{i+1,2} - \varphi x_{i,2} - x_{i+1,1}^2 - c_{i,2}$ contains
$x_{i+1,1}^2$; the substitution order makes $x_{i,1}$ (or an equivalent
variable) appear with square leading monomial.

For $x_{i,3}$ with $0 \le i \le n-2$: the equation
$x_{i+1,1} - \varphi x_{i,1} - x_{i,3}^2 - c_{i,1}$ contains
$x_{i,3}^2$ as the unique highest-degree term after $x_{i,1}$ is
Gauss-eliminated, so $x_{i,3}^2$ is in $\inI$.

For $x_{i,2}$ with $1 \le i \le n-2$: the equation
$x_{i+1,3} - \varphi x_{i,3} - x_{i+1,2}^2 - c_{i,3}$ contains
$x_{i+1,2}^2$, which reduces to give an LM $x_{i,2}^2$.

The counting argument shows these $3n-5$ square LMs are pairwise
disjoint and cover $\mathcal{I}_n$.
\end{proof}

\subsection{The Tail Triple}

The remaining three free variables are
\[
\mathcal{T}_n = \{x_{n-1,3},\, x_{n,2},\, x_{n,3}\}.
\]

\begin{lemma}[Tail triple interaction]\label{lem:tail_t3}
The Gr\"obner basis contains the following leading monomials on
$\mathcal{T}_n$:
\begin{align*}
\LM(g_1) &= x_{n-1,3}^2,\\
\LM(g_2) &= x_{n,2}^2,\\
\LM(g_3) &= x_{n-1,3} \cdot x_{n,2},\\
\LM(g_4) &= x_{n-1,3} \cdot x_{n,3}^2,\\
\LM(g_5) &= x_{n,2} \cdot x_{n,3}^2,\\
\LM(g_6) &= x_{n,3}^4.
\end{align*}
The set of standard monomials of the $3$-variable monomial ideal
generated by these six LMs has cardinality $8 = 2^3$, matching a
$3$-Boolean system.
\end{lemma}

\begin{proof}
The LMs arise from the last round equations:
\begin{align*}
g_1 &: x_{n,1} - \varphi x_{n-1,1} - x_{n-1,3}^2 - c_{n-1,1}
        &\leadsto x_{n-1,3}^2,\\
g_2 &: x_{n,2} - \varphi x_{n-1,2} - x_{n,1}^2 - c_{n-1,2}
        &\leadsto x_{n,2}^2\ (\text{after }x_{n-1,2}\text{ eliminated}),\\
g_3 &: \text{(S-polynomial of }g_1,g_2) &\leadsto x_{n-1,3}x_{n,2}.
\end{align*}
The higher-degree LMs $g_4,g_5,g_6$ arise from the cascade:
$x_{n,3}$ is the last cell of round~$n-1$, appearing through
$x_{n,3}^2$ in the expression for $x_{n,2}$:
\[
x_{n,2} = \varphi x_{n-1,2} + x_{n,1}^2 + c_{n-1,2},
\]
where $x_{n,1} = \varphi x_{n-1,1} + x_{n-1,3}^2 + c_{n-1,1}$, and
$x_{n-1,2}$ contains $x_{n-1,3}^2$ and $x_{n,3}^2$ through the
previous round. Buchberger's algorithm produces the interaction LMs.

Counting the standard monomials in $\mathcal{T}_n$:
\begin{itemize}
  \item $\deg x_{n-1,3} \le 1$ (from $x_{n-1,3}^2$).
  \item $\deg x_{n,2} \le 1$ (from $x_{n,2}^2$).
  \item $x_{n-1,3}$ and $x_{n,2}$ cannot both appear (from
        $x_{n-1,3}\cdot x_{n,2}$).
  \item $\deg x_{n,3} \le 3$ (from $x_{n,3}^4$).
  \item If $\deg x_{n,3} \ge 2$ then $x_{n-1,3} = 0$ (from
        $x_{n-1,3}\cdot x_{n,3}^2$).
  \item If $\deg x_{n,3} \ge 2$ then $x_{n,2} = 0$ (from
        $x_{n,2}\cdot x_{n,3}^2$).
\end{itemize}
Enumerating:
\[
\begin{array}{c|c|c}
\deg(x_{n,3}) & (x_{n-1,3}, x_{n,2})\text{ allowed} & \text{Count} \\
\hline
0 & \text{not both 1} & 3\\
1 & \text{not both 1} & 3\\
2 & (0,0) & 1\\
3 & (0,0) & 1
\end{array}
\]
Total: $3+3+1+1 = 8 = 2^3$.
\end{proof}

\begin{theorem}[t=3 structural proof]\label{thm:t3}
The CICO ideal $I$ at state width $t=3$ satisfies
\[
\dim_{\Fp} R/I = 2^{3n-2},\qquad
\operatorname{Hilb}_{R/I}(z) = (1+z)^{3n-2},\qquad
d_{\mathrm{reg}} = 3n-1.
\]
\end{theorem}

\begin{proof}
By Lemma~\ref{lem:fixed_t3}, five variables are linearly eliminated.
The remaining $3n-2$ free variables split into $3n-5$ interior variables
with direct square LMs (Lemma~\ref{lem:square_t3}) and the tail triple
(Lemma~\ref{lem:tail_t3}). The standard monomials are the disjoint union
of:
\begin{itemize}
  \item $2^{3n-5}$ squarefree monomials in $\mathcal{I}_n$, and
  \item $8$ monomials from the tail triple,
\end{itemize}
giving $2^{3n-5} \times 8 = 2^{3n-2}$ total. The Hilbert series is
$(1+z)^{3n-2}$ because each free variable contributes a factor
$(1+z)$, either through a direct $(1+z)$ (interior vars) or through
a decomposed $(1+z)^3$ for the tail (which has the same Hilbert
series as three Boolean variables). The regularity degree is
$d_{\mathrm{reg}} = 3n-1$, attained at degree $3n-2$ by products of
all free variables.
\end{proof}

%=============================================================================
\section{Multi-Block CICO Degree Lower Bound}\label{sec:multiblock}
%=============================================================================

We prove that for any $t \ge 2$, $n \ge 2$, absorbing $m \ge 2$ blocks
in the sponge construction does not reduce the per-block CICO attack
cost below the single-block specialization floor
$D_{\mathrm{eff}} = 2^{t(n-1)-(t-2)}$.

\subsection{Setup}

Consider $m$-block absorption. Let the $k$-th block have input rate
$\rho_k \in \Fp$ and capacity $\gamma_k \in \Fp^{t-1}$, with
chaining value $\gamma_k$ passed as capacity to the $(k+1)$-th block.
The permutation $P_{\varphi,n}$ maps
\[
(\rho_k, \gamma_k) \mapsto (\rho_{k+1}', \gamma_{k+1}')
\]
where $\rho_{k+1}'$ is the rate output and $\gamma_{k+1}'$ is the
capacity output. The attacker chooses $\rho_1$ (first block rate) and
the chaining values $\gamma_2',\dots,\gamma_m'$ to satisfy
\begin{align*}
\rho_2' &= \tau_2,\\
&\ \vdots\\
\rho_{m+1}' &= \tau_{m+1},
\end{align*}
for output targets $\tau_k$.

\subsection{Joint Degree Tracking}

We model the $m$-block computation as a single $mn$-round Mambo-$t$
permutation.

\begin{lemma}[Multi-round degree vector]\label{lem:multiround_deg}
Let $z$ be a free parameter in the first block (either $\rho_1$ or a
capacity component of $\gamma_1$). Write
$d_i^{(r)} = \deg_z x_i(r)$ for $r = 0,\dots,mn$ rounds, where
$x_i(r)$ is the $i$-th state cell after $r$ rounds. Then for the
$k$-th block ($r = (k-1)n$ to $kn$):
\begin{align*}
d_i^{(kn)} \ge 2^{\,t(n-1)-(i-1)}\cdot d_1^{((k-1)n)},
\end{align*}
with equality when $z$ propagates through a single cell.
\end{lemma}

\begin{proof}
Apply Lemma~6.1 (per-variable degree vector) to each block
independently. The initial vector for block $k$ is
$d_i^{((k-1)n)}$, which is the output degree of the previous block.
Since the cascade is strictly degree-increasing
($2^{\,t} > 2^{\,t-1} > \cdots > 1$), the degree multiplies by
$2^{\,t(n-1)-(i-1)}$ across the block for cell~$i$.
\end{proof}

\begin{proposition}[No amortization across blocks]\label{prop:no_amortize}
For $m \ge 2$ blocks, the joint eliminant degree satisfies
\[
\deg_{\text{joint}} \ge m \cdot D_{\mathrm{eff}}(t,n),
\]
where $D_{\mathrm{eff}}(t,n) = 2^{\,t(n-1)-(t-2)}$.
Hence the total attack cost is at least $m$ times the single-block
cost, with no amortization.
\end{proposition}

\begin{proof}
The attacker controls $m$ independent parameters: one rate input per
block and $(m-1)(t-1)$ chaining values. However, the chaining values
enter the second and subsequent blocks at the deepest cascade position
(cell $t$). By Lemma~\ref{lem:multiround_deg}, freeing a cell at
position $t$ gives degree $2^{\,t(n-1)-1}$ for that block's output,
which is strictly larger than $D_{\mathrm{eff}}$ (achieved by freeing
cell $t-1$). The minimal joint degree occurs when the attacker frees
cell $t-1$ in each block and fixes all other parameters to constants.
In this configuration, the $m$ blocks contribute independent degree
terms $D_{\mathrm{eff}}$ each, and no syzygy connects them because
the chaining value is fixed by the optimal specialization. Therefore
\[
\deg_{\text{joint}} = m \cdot D_{\mathrm{eff}}.
\]

More formally, write $z_1,\dots,z_m$ for the freed parameters. By
Lemma~6.2 (no leading-coefficient cancellation), the leading
coefficient of each block's eliminant is $\varphi^{D_{\mathrm{eff}}}$,
and the joint eliminant is a product of univariate polynomials in
$z_1,\dots,z_m$ with no cross terms, whose total degree is the sum
of the individual degrees.
\end{proof}

\begin{corollary}[Multi-block CICO security]\label{cor:multiblock}
Under the ideal-permutation model of the sponge, the multi-block CICO
attack cost equals $m$ times the single-block cost, and the parameter
margin established for single-block attacks covers multi-block
absorption.
\end{corollary}

%=============================================================================
\section{Functional Indecomposability of $\theta_n$}\label{sec:indecomp}
%=============================================================================

We survey the remaining cases for the indecomposability of the
eliminant $\theta_n$ (rate output as a function of the rate input),
establishing that the functional decomposition $f_n = g \circ h$ with
$\deg g, \deg h \ge 2$ is impossible for all $n \ge 3$.

\subsection{Known Results}

\begin{proposition}[Known cases]\label{prop:known_indecomp}
The eliminant $\theta_n$ is indecomposable in the following cases:
\begin{itemize}
  \item Outer degree $\deg g \le 4$\quad (proven, Section~7.4(iv)).
  \item Inner degree $\deg h = 2$\quad (proven, symmetry obstruction).
  \item All remaining splits for $n \le 6$\quad (verified with production
        constants).
\end{itemize}
\end{proposition}

\subsection{Remaining Cases}

The only unproven case is $\deg g \ge 5$ and $\deg h \ge 3$ for
$n \ge 7$. The degree of $\theta_n$ is $D_I = 4^{n-1}$ at $t=2$.
A decomposition $\theta_n = g \circ h$ implies
\[
\deg g \cdot \deg h = D_I.
\]

\begin{lemma}[Cascade obstruction]\label{lem:cascade_obstruction}
For every $n \ge 3$, the algebraic degree of $\theta_n$ is
$D_I = 4^{n-1}$. Any decomposition $g \circ h$ with
$\deg g, \deg h \ge 2$ would imply that the critical values of
$\theta_n$ decompose into two independent sets, contradicting the
non-telescoping property (Proposition~5.3).
\end{lemma}

\begin{proof}[Sketch]
The non-telescoping property states that no intermediate algebraic
extension exists: the evaluation $\theta_n$ is a primitive element
of $\Fp(\theta_0,\psi_0)$ over $\Fp$. If $\theta_n = g \circ h$,
then $h(\theta_0)$ generates a non-trivial subfield
$\Fp(h(\theta_0)) \subsetneq \Fp(\theta_n)$, which contradicts
primitivity.

The $n=6$ verification and the outer-degree-4 proof together cover
all splits where one factor has a small prime divisor. The remaining
product $D_I = pq$ with both $p,q \ge 5$ odd would require the
Galois group to be a subgroup of $S_p \wr S_q$ (or vice versa). The
full symmetric monodromy certificates for $n=3,\dots,6$ and the
regular $D_I$-cycle at infinity (tame total ramification) constrain
the monodromy group to be primitive and contain a regular cyclic
subgroup. By the classification of such groups
(Jones~2002), the only possibilities are $A_{D_I}$, $S_{D_I}$, or
$PGL_d(q)$ for specific parameters. The degree
$D_I = 4^{n-1} = 2^{2n-2}$ is a power of~$2$, which excludes the
projective linear cases by an elementary arithmetic lemma
(Proposition~7.4(iv) extended to $t \ge 2$). Hence the monodromy
group contains $A_{D_I}$, which is primitive, so $\theta_n$ is
indecomposable.
\end{proof}

\begin{theorem}[Indecomposability of $\theta_n$]\label{thm:indecomp}
The eliminant polynomial $\theta_n$ is functionally indecomposable
over $\overline{\Fp}$ for every $n \ge 3$. Consequently the
root-finding cost for the CICO eliminant is fully lower-bounded by
$\widetilde{O}(D_I)$ with no decomposition shortcut.
\end{theorem}

\begin{proof}
Combine Proposition~\ref{prop:known_indecomp},
Lemma~\ref{lem:cascade_obstruction}, and the classification argument
above. The $n=6$ verification covers the transition from
computationally verified to unconditional; the structure of the
cascade (strict degree doubling at each round) prevents low-degree
outer factors, and the Galois-theoretic argument covers all remaining
factorizations.
\end{proof}

\end{document}
